\chapter{Benchmark Description}
\label{app:benchmark}

This appendix provides a detailed account of the hard synthetic benchmark
used in the experimental evaluation of Chapter~\ref{ch:experiments}.  The
benchmark was designed to overcome three limitations of ad hoc
Agile-SOFL example tasks: insufficient scale for statistical comparison,
insufficient difficulty to discriminate among pipeline modes, and
insufficient control over the structural properties (guard overlap, scenario
count, output arity) that the HSP-Agile pipeline is intended to enforce.
Published Agile-SOFL examples~\citep{liu2014sofl,li2022agilesofl} provide seed tasks but lack
the scale and difficulty control required for mode discrimination.

% ---------------------------------------------------------------------------
\section{Design Rationale and Task Selection Criteria}
\label{sec:bench:criteria}
% ---------------------------------------------------------------------------

Each benchmark task must satisfy four admissibility criteria before inclusion
in the evaluation corpus:

\paragraph{C1 (Satisfiability).}
Every FSF scenario, including the \texttt{others} default, must be
individually satisfiable: the conjunction of the scenario guard and
its output postcondition must admit at least one model in the input
domain.  Tasks failing this check are discarded.

\paragraph{C2 (Implementability).}
At least one correct implementation must exist that evaluates guards
in specification order and produces the prescribed outputs on every
scenario.  An auto-generated reference implementation serves as a
constructive witness.

\paragraph{C3 (Discriminative difficulty).}
Guard regions must overlap: at least two scenario preconditions must
be simultaneously satisfiable on a non-empty input subspace, so that
incorrect guard ordering produces detectable errors on boundary inputs
rather than being masked by accidental correctness.

\paragraph{C4 (Uniform structure).}
All hard tasks share an identical signature shape (five natural-number
inputs, three natural-number outputs, eight ordered scenarios) so that
cross-task comparisons reflect algorithmic differences rather than
incidental variation in task complexity.

The first twenty tasks in the broader corpus were extracted manually from
published Agile-SOFL example specifications~\citep{liu2014sofl,liu2004sofl} and retained as a development set.
They were excluded from the primary evaluation to prevent information leakage
between prompt-engineering iterations and held-out testing.  The hard synthetic
suite described below constitutes the 120-task evaluation benchmark.

% ---------------------------------------------------------------------------
\section{Hard Benchmark Task Categories}
\label{sec:bench:categories}
% ---------------------------------------------------------------------------

Hard tasks are produced by a programmatic generator that parameterises task
difficulty through two structural mechanisms: \emph{precedence-sensitive
guard construction} and \emph{multi-output coupling}.

\subsection{Signature and scenario template}

Every generated task defines a process with the following fixed signature:

\begin{itemize}
  \item \textbf{Inputs} ($|\mathbf{x}| = 5$): natural-number variables
        $a, b, c, d, e$, each ranging over $[0, 100]$.
  \item \textbf{Outputs} ($|\mathbf{y}| = 3$): natural-number variables
        \textit{risk}, \textit{action}, and \textit{score}, representing a
        composite decision tuple (risk tier, recommended action, aggregate
        score).
  \item \textbf{Scenarios} ($|\mathcal{S}_t| = 8$): seven guarded scenarios
        plus one \texttt{others} default clause, evaluated in listed order
        with top-down precedence.
\end{itemize}

The uniform arity ensures that formal conformance checking, pattern guard
analysis, and mutation testing operate under consistent structural assumptions
across all 120 evaluation tasks.

\subsection{Precedence-sensitive guards}

Each guarded scenario's pre-condition is a conjunction of two or three
relational predicates over the input variables, with thresholds drawn from
task-specific integer constants.  The generator maintains a pool of nine
candidate guard--postcondition pairs spanning diverse input combinations
(high-$a$/low-$b$ regions, zero-input boundaries, multi-variable conjunctions).
For each task, a subset of seven pairs is selected and shuffled into priority
order; the \texttt{others} scenario closes the specification with a zero
default output tuple.

Critically, the guard templates are designed so that multiple scenarios'
preconditions overlap in the input space.  Correct behaviour therefore
requires strict priority evaluation: an implementation that tests guards in
reverse order or that collapses overlapping regions into a single branch will
produce correct outputs on most inputs yet fail on witnesses lying in the
overlap region.  This property directly targets scenario-order violation
(SOV), the most prevalent fault class in LLM-generated FSF implementations~\citep{dakhel2023copilot,chen2021codex}.

\subsection{Multi-output coupling}

Each scenario's postcondition binds all three output variables simultaneously
(e.g.\ \code{risk eq 3 \&\& action eq 4 \&\& score eq a}), preventing
implementations from optimising for a single output field while neglecting
others.  Output expressions may reference input variables, constants, or
combinations thereof, requiring the implementation to propagate input values
into structured return dictionaries rather than returning scalar results.

% ---------------------------------------------------------------------------
\section{Deterministic Generation and Validation}
\label{sec:bench:validation}
% ---------------------------------------------------------------------------

Task generation is fully deterministic: a fixed pseudorandom seed (42) governs
threshold sampling, scenario subset selection, and guard shuffling.  Re-running
the generator with the same seed reproduces the identical 180-candidate task
pool.

Each candidate task passes through a two-stage validation pipeline before
admission to the evaluation corpus:

\begin{enumerate}
  \item \textbf{Satisfiability check.}  For each of the eight scenarios, an
        SMT solver verifies that the guard conjunction combined with output
        postconditions is satisfiable over the declared input domain~\citep{demoura2008z3,barrett2018smt}.  Tasks
        with even one unsatisfiable scenario are rejected.

  \item \textbf{Reference implementation generation.}  A correct reference
        implementation is synthesised directly from the FSF specification by
        translating each scenario into an ordered \code{if}--\code{elif}
        branch with a terminal \code{else} for the \texttt{others} clause.
        The reference serves as ground truth for mutation testing and as a
        sanity check that the specification is implementable.
\end{enumerate}

Of 180 generated candidates, 60 failed the satisfiability filter, yielding
the final benchmark of 120 validated hard tasks.  Table~\ref{tab:bench:params}
(below) summarises the generation parameters.

\begin{table}[htbp]
  \caption{Hard benchmark generation parameters.}
  \label{tab:bench:params}
  \centering
  \begin{tabular}{ll}
    \toprule
    Parameter & Value \\
    \midrule
    Input variables       & 5 (\code{a}, \code{b}, \code{c}, \code{d}, \code{e}) \\
    Output variables      & 3 (\code{risk}, \code{action}, \code{score}) \\
    Scenarios per task    & 8 (7 guarded + 1 \texttt{others}) \\
    Input domain          & $[0, 100]^5$ (natural numbers) \\
    Guard overlap         & By construction (multi-scenario overlap) \\
    Random seed           & 42 \\
    Candidates generated  & 180 \\
    Candidates retained   & 120 \\
    \bottomrule
  \end{tabular}
\end{table}

% ---------------------------------------------------------------------------
\section{FSF Specification Format}
\label{sec:bench:format}
% ---------------------------------------------------------------------------

Each task is encoded as a structured FSF artefact comprising a process
signature, an ordered scenario list, and a natural-language prompt
specification.  The on-disk representation uses the following fields:

\begin{description}
  \item[\code{signature}.]
        Declares input and output variables with type annotations
        (\code{nat} for natural numbers).
  \item[\code{fsfScenarios}.]
        An ordered list of scenario records.  Each guarded scenario carries
        a \code{test} field (the guard predicate in Agile-SOFL relational
        notation, using \code{gt}, \code{le}, \code{eq}, and \code{\&\&}) and
        a \code{def} field (the output postcondition).  The final scenario
        has \code{kind = others} with \code{test = "others"}.
  \item[\code{promptSpec}.]
        A human-readable rendering of the FSF specification, including an
        explicit instruction to evaluate conditions in listed order
        (top-down precedence).
  \item[\code{referenceCode}.]
        Auto-generated Python implementation satisfying the specification
        under ordered guard semantics.
\end{description}

The prompt specification deliberately states precedence ordering because
prior work shows that LLMs frequently generate \code{if}--\code{elif} chains
without respecting specification-level priority when this instruction is
omitted~\citep{lynn2024verifierloop,loughman2024dafnybench}.  Including it isolates the effect of formal checking and pattern
guarding from pure prompt-compliance effects.

% ---------------------------------------------------------------------------
\section{Illustrative Task Structure}
\label{sec:bench:example}
% ---------------------------------------------------------------------------

Although individual guard thresholds vary across tasks, the structural template
is uniform.  A representative task (with simplified threshold values) reads:

\begin{quote}\small\ttfamily
Process HardCase001(a,b,c,d,e) -> (risk,action,score)\\[2pt]
FSF specification:\\
if (a gt 5 \&\& b gt 7 \&\& c gt 4) => risk eq 3 \&\& action eq 4 \&\& score eq a\\
if (a gt 5 \&\& b gt 7 \&\& c le 4) => risk eq 2 \&\& action eq 3 \&\& score eq b\\
if (a gt 5 \&\& b le 7 \&\& d gt 3) => risk eq 2 \&\& action eq 2 \&\& score eq d\\
\ldots\\
if (a eq 0 \&\& b eq 0) => risk eq 0 \&\& action eq 5 \&\& score eq 0\\
others => risk eq 0 \&\& action eq 0 \&\& score eq 0\\[2pt]
Important: evaluate conditions in listed order (top-down precedence).
\end{quote}

The first two scenarios share the prefix \code{a gt 5 \&\& b gt 7} but
differ on \code{c}; any input satisfying the first scenario's guard also
satisfies the second's guard prefix.  An implementation that evaluates the
second scenario before the first will assign incorrect outputs to inputs in
the first scenario's exclusive region---a violation detectable by
SMT-generated boundary witnesses.

% ---------------------------------------------------------------------------
\section{Mutation Operators and Evaluation Protocols}
\label{sec:bench:mutation}
% ---------------------------------------------------------------------------

The benchmark supports two complementary evaluation protocols beyond the
primary mode comparison:

\paragraph{Specification-confusion mutants.}
Four operators inject faults at the specification layer: predicate negation
(SNO), relational-operator swap (ORO), scenario deletion (MCO), and boundary
weakening (BCO)~\citep{demillo1978mutation,offutt1992mutation}.  The prevention protocol applies each operator to the FSF
specification while holding the reference implementation fixed, then tests
whether the formal checker rejects the now-mismatched pair.

\paragraph{Implementation-screening mutants.}
Four operators inject faults at the implementation layer: condition inversion
(ICO), wrong primary output (WRO), off-by-one comparison (MBO), and else-branch
removal (DRO).  The prevention protocol applies each operator to the reference
implementation and tests whether the hybrid acceptance predicate
($\mathrm{Conf} = 1 \wedge \mathcal{P}^H = \emptyset$) rejects the mutant.

Together, the eight mutation operators provide coverage of both specification
misinterpretation and implementation-level coding defects, aligned with the
pattern catalogue of Appendix~\ref{app:patterns}.

% ---------------------------------------------------------------------------
\section{Benchmark Statistics}
\label{sec:bench:stats}
% ---------------------------------------------------------------------------

Table~\ref{tab:bench:summary} (below) summarises the structural properties of the
120-task evaluation benchmark.

\begin{table}[htbp]
  \caption{Summary statistics of the 120-task hard benchmark.}
  \label{tab:bench:summary}
  \centering
  \begin{tabular}{lc}
    \toprule
    Statistic & Value \\
    \midrule
    Tasks in evaluation corpus        & 120 \\
    Inputs per task                     & 5 \\
    Outputs per task                    & 3 \\
    Scenarios per task                  & 8 \\
    Guarded scenarios                   & 7 \\
    Default (\texttt{others}) scenarios & 1 \\
    Mean guard overlap rate             & 1.8 scenarios/input \\
    Mean SMT models per scenario        & 312 (cap 1000) \\
    Total experimental runs (7 modes)   & 840 \\
    \bottomrule
  \end{tabular}
\end{table}

The mean guard overlap rate of 1.8 indicates that a randomly sampled input
satisfies approximately 1.8 scenario guards on average, requiring correct
priority resolution.  The mean of 312 satisfying models per scenario confirms
that no scenario is degenerate or vacuously satisfiable only at a single
point.  To our knowledge, this benchmark constitutes the largest published
collection of FSF tasks with SMT-certified scenario validity and controlled
precedence-sensitive guard structure.
